% Einheit 1 von 3 – Entscheidungsregel bestimmen: die Nullhypothese liefert p und die Binomialverteilung der Testgröße, die Alternative die Seite des Ablehnungsbereichs (rechtsseitig bei „mehr als“, linksseitig bei „weniger als“); die Grenze über kumulierte Wahrscheinlichkeiten so wählen, dass das Signifikanzniveau eingehalten wird (bank/hypothesentests/e1.jsonl, zusammenbau v0.1)
%% TODO Zeitmarke Einheit 1: Marken-Zeile nicht lesbar
%% TODO Zweigzeile Teil 1: was hier gelernt wird (Satz in Schülersprache aus dem Einheitstitel) – Einheit 1
\einheitenkopf[e1]{Einheit 1 von 3 · Entscheidungsregel bestimmen: die Nullhypothese liefert p und die Binomialverteilung der Testgröße, die Alternative die Seite des Ablehnungsbereichs (rechtsseitig bei „mehr als“, linksseitig bei „weniger als“); die Grenze über kumulierte Wahrscheinlichkeiten so wählen, dass das Signifikanzniveau eingehalten wird}
\zweigzeile{Abitur LK}
\setcounter{aufgabe}{7}

%% TODO Titel: Ich-kann-Satz fehlt in der Bank (Platzhalter: Kettenname) – Nr. 8
%% TODO Anweisung: ein Satz über der ersten Teilaufgabe fehlt in der Bank – Nr. 8
\begin{aufgabe}{Entscheidungsregel}
%% TODO \rechenplatz{n}: Schrittzahl des Grundfalls fehlt in der Bank (nur bei fester Schreibform, 2.3 b)
\begin{teile}
\teil $H_0\colon p \le 0{,}2$, Gegenhypothese $H_1\colon p > 0{,}2$. Auf welcher Seite liegt der Ablehnungsbereich? \\ \kreuz{rechtsseitig} \\ \kreuz{linksseitig}
\teil Eine Buchhändlerin vermutet, dass mehr als 25\,\% der Jugendlichen der Stadt Kunden in ihrem Laden sind; sonst plant sie eine Werbeaktion. Um unnötige Kosten zu vermeiden, soll die Nullhypothese $H_0$: „Höchstens 25\,\% der Jugendlichen sind Kunden“ mit 120 zufällig befragten Jugendlichen auf dem Signifikanzniveau 5\,\% getestet werden; $X$ ist die Anzahl der Kunden unter ihnen. Bestimme die Entscheidungsregel. \hfill (Abitur 2018 LK)
\teil Ein Reiseveranstalter verlängert ein Gewinnspiel nur, wenn mindestens 4\,\% der Teilnehmer danach eine Reise buchen. Getestet wird $H_0\colon p \ge 0{,}04$ ($p$ Buchungswahrscheinlichkeit) auf dem Signifikanzniveau 5\,\% mit 650 Teilnehmern; $X$ ist die Anzahl der Buchungen. Bestimme die Entscheidungsregel. \hfill (Abitur 2023 LK)
\teil Ein Fährbetrieb setzt eine zusätzliche Fähre für Radfahrer nur fort, wenn $H_0$: „Höchstens 12\,\% der Fahrgäste haben ein Rad dabei“ abgelehnt wird. Getestet wird auf dem Signifikanzniveau 8\,\% an 400 Fahrgästen; $X$ ist die Anzahl der Fahrgäste mit Rad. Bestimme die Entscheidungsregel und sichere die Grenze am Nachbarwert ab. \hfill (Abitur 2025 LK)
\teil Ein Streamingdienst testet $H_0$: „Höchstens 40\,\% der Nutzer sehen eine Serie zu Ende“ an 150 Nutzern auf dem Signifikanzniveau 5\,\%. Als Ablehnungsbereich wird $\{71; \ldots; 150\}$ angegeben, mit den Lösungsschritten: „$Y$ Anzahl der Nutzer, die die Serie zu Ende sehen; $P(Y \ge 71) \approx 0{,}041$.“ Begründe, warum diese Schritte nicht ausreichen, und ergänze sie. \hfill (Abitur 2024 LK)
\end{teile}
\end{aufgabe}

%% TODO Titel: Ich-kann-Satz fehlt in der Bank (Platzhalter: Kettenname) – Nr. 9
%% TODO Anweisung: ein Satz über der ersten Teilaufgabe fehlt in der Bank – Nr. 9
\begin{aufgabe}{Ablehnungsgrenze bei größerem Stichprobenumfang ohne höheren Fehler erster Art bestimmen}
\begin{teile}
\teil Ein Hersteller will höchstens 15\,\% fehlerhafte Teile. Bisher prüft er 80 Teile und bessert das Verfahren nach, wenn mindestens 18 fehlerhaft sind. Künftig prüft er 160 Teile; die Wahrscheinlichkeit für eine unnötige Nachbesserung soll dadurch nicht steigen. Ab wie vielen fehlerhaften Teilen wird nun nachgebessert? \hfill (Abitur 2017 LK)
\end{teile}
\end{aufgabe}

%% TODO Titel: Ich-kann-Satz fehlt in der Bank (Platzhalter: Kettenname) – Nr. 10
%% TODO Anweisung: ein Satz über der ersten Teilaufgabe fehlt in der Bank – Nr. 10
\begin{aufgabe}{Entscheidungsregel – Fehler finden, Begründen, Anwendung}
\begin{teile}
\teil Getestet wird $H_0$: „Mindestens 65\,\% der Kunden sind zufrieden“ an 90 Kunden auf dem Signifikanzniveau 5\,\%. Lisa bestimmt das kleinste $k$ mit $P(X \ge k) \le 0{,}05$ und erhält $k = 67$. Finde den Fehler.
\teil Begründe, warum bei $H_0\colon p \le 0{,}35$ gegen $H_1\colon p > 0{,}35$ der Ablehnungsbereich rechts liegt.
\teil Ein Kino kauft eine zweite Popcornmaschine nur, wenn $H_0$: „Höchstens 40\,\% der Besucher kaufen Popcorn“ abgelehnt wird; die Regel für 60 befragte Besucher lautet: Ablehnung ab 31 Käufern. In der Stichprobe kaufen 29. Was folgt daraus?
\end{teile}
\end{aufgabe}
